\documentclass[11pt, a4paper]{article}

\usepackage[a4paper, margin=2.5cm]{geometry}
\usepackage{fontspec}
\usepackage[russian, bidi=basic, provide=*]{babel}
\babelprovide[import, onchar=ids fonts]{russian}
\babelprovide[import, onchar=ids fonts]{english}
\babelfont{rm}{Noto Sans}
\babelfont{sf}{Noto Sans}

\usepackage{graphicx}
\usepackage{float}     
\usepackage{xcolor}
\usepackage{listings}
\usepackage{enumitem}
\usepackage{amsmath}
\usepackage{booktabs}
\usepackage{tabularx}
\usepackage{titlesec}
\usepackage{hyperref}

\setlist[itemize]{label=-, itemsep=3pt, parsep=0pt, topsep=5pt}
\setlist[enumerate]{itemsep=3pt, parsep=0pt, topsep=5pt}
\setlength{\parskip}{0.3em}
\setlength{\parindent}{1em}
\titlespacing*{\section}{0pt}{1em}{0.5em}
\titlespacing*{\subsection}{0pt}{0.8em}{0.4em}

% === Настройка подсветки кода ===
\lstset{
    language=C++,
    basicstyle=\ttfamily\small,
    keywordstyle=\color{blue}\bfseries,
    commentstyle=\color{gray}\itshape,
    stringstyle=\color{red},
    numbers=left,
    numberstyle=\tiny\color{gray},
    frame=single,
    breaklines=true,
    captionpos=b,
    tabsize=4,
    showstringspaces=false
}

\hypersetup{
    colorlinks=true,
    linkcolor=black,
    urlcolor=blue
}

\begin{document}

\begin{titlepage}
    \begin{center}
        \Large\textbf{МИНИСТЕРСТВО НАУКИ И ВЫСШЕГО ОБРАЗОВАНИЯ РОССИЙСКОЙ ФЕДЕРАЦИИ} \\
        \vspace{0.3cm}
        \large Федеральное государственное автономное образовательное учреждение высшего образования \\
        \textbf{«Московский авиационный институт (национальный исследовательский университет)»}
    \end{center}

    \vspace{3cm}

    \begin{center}
        \huge\textbf{ОТЧЁТ ПО ЛАБОРАТОРНОЙ РАБОТЕ № 2}
    \end{center}

    \vspace{1.5cm}

    \begin{center}
        \large\textbf{Дисциплина:} Операционные системы и архитектура компьютера \\
        \vspace{0.5cm}
        \large\textbf{Тема:} Управление процессами и межпроцессное взаимодействие через каналы (pipe)
    \end{center}

    \vspace{3cm}

    \begin{flushright}
        \begin{tabular}{ll}
            \textbf{Выполнил:} & студент группы М8О-203БВ-25 \\
                               & Воскресенский И.А. \\[1cm]
            \textbf{Проверил:} & Преподаватель \\
                               & Миронов Е.С. \\
        \end{tabular}
    \end{flushright}

    \vfill

    \begin{center}
        \large Москва, 2026 г.
    \end{center}
\end{titlepage}

\section{Цель работы}
Приобретение практических навыков в управлении процессами в операционной системе и обеспечении обмена данными между процессами посредством каналов (pipe). В ходе работы необходимо разработать программу на языке C++, в которой родительский процесс создаёт дочерний, передаёт ему данные через неименованный канал, а дочерний процесс выполняет вычисления и записывает результат в файл.

\section{Краткое описание архитектуры решения}
Решение состоит из двух независимых исполняемых модулей: \texttt{parent} (родительский процесс) и \texttt{child} (дочерний процесс), взаимодействующих через неименованный канал POSIX (\texttt{pipe}).

\textbf{Родительский процесс} выполняет следующие действия:
\begin{enumerate}
    \item запрашивает у пользователя имя выходного файла;
    \item создаёт канал функцией \texttt{pipe()};
    \item порождает дочерний процесс через \texttt{fork()};
    \item в дочерней ветке перенаправляет стандартный ввод на читающий конец канала через \texttt{dup2()} и запускает программу \texttt{child} с помощью \texttt{execl()};
    \item в родительской ветке читает ввод пользователя и записывает его в пишущий конец канала;
    \item ожидает завершения дочернего процесса через \texttt{waitpid()}.
\end{enumerate}

\textbf{Дочерний процесс}:
\begin{enumerate}
    \item принимает имя файла из аргументов командной строки;
    \item открывает файл для записи;
    \item в цикле читает строки из \texttt{stdin} (связанного с каналом);
    \item парсит числа типа \texttt{float} и вычисляет их сумму;
    \item записывает результат в файл.
\end{enumerate}

\subsection{Схема взаимодействия компонентов}
Поток данных и управляющих сообщений между компонентами:
\begin{itemize}
    \item \texttt{User} $\xrightarrow{\text{stdin}}$ \texttt{Parent (PID P)}
    \item \texttt{Parent} $\xrightarrow{\text{write(pipe\_fd[1])}}$ \textbf{pipe} $\xrightarrow{\text{read(STDIN)}}$ \texttt{Child (PID C)}
    \item \texttt{Child} $\xrightarrow{\text{ofstream}}$ \texttt{result.txt}
\end{itemize}
Управляющее сообщение: закрытие пишущего конца канала родителем генерирует EOF для дочернего процесса, сигнализируя о завершении ввода.

\section{Листинг ключевых фрагментов кода}
Полные исходные файлы доступны в репозитории: \texttt{https://github.com/friceXXX/osi-lab2}

\subsection{Создание канала и порождение процесса}
Функция \texttt{pipe()} создаёт пару связанных файловых дескрипторов. После \texttt{fork()} каждый процесс закрывает неиспользуемый конец канала, что необходимо для корректной генерации EOF.

\begin{lstlisting}[caption={Создание канала и fork}, label={lst:pipe}]
if (pipe(pipefd) == -1) {
    ERR_EXIT("pipe");
}

cpid = fork();
if (cpid == -1) {
    ERR_EXIT("fork");
}

if (cpid == 0) {
    if (close(pipefd[1]) == -1) ERR_EXIT("close write-end");
    if (dup2(pipefd[0], STDIN_FILENO) == -1) ERR_EXIT("dup2");
    if (close(pipefd[0]) == -1) ERR_EXIT("close old read-end");

    execl("./child", "child", filename.c_str(), (char *)NULL);
    ERR_EXIT("execl");
} else {
    if (close(pipefd[0]) == -1) ERR_EXIT("close read-end");
    if (close(pipefd[1]) == -1) ERR_EXIT("close write-end");
    int status;
    if (waitpid(cpid, &status, 0) == -1) {
        ERR_EXIT("waitpid");
    }
}
\end{lstlisting}

\subsection{Парсинг чисел и вычисление суммы}
Дочерний процесс использует \texttt{std::istringstream} для извлечения чисел типа \texttt{float} из строки. Цикл автоматически завершается при исчерпании потока.

\begin{lstlisting}[caption={Парсинг float и подсчет суммы}, label={lst:parse}]
while (std::getline(std::cin, line)) {
    std::istringstream iss(line);
    float value = 0.0f, sum = 0.0f;
    int count = 0;
    
    while (iss >> value) {
        sum += value;
        ++count;
    }
    
    if (count > 0) {
        out_file << "Sum of " << count 
                 << " numbers: " << sum << "\n";
    } else ...
    out_file.flush();
}
\end{lstlisting}

\section{Результаты работы}

\subsection{Консольный вывод или логи}
Ниже приведён фактический вывод программы при тестировании. Различные PID родителя и дочернего процесса подтверждают конкурентное выполнение двух независимых процессов.

\begin{verbatim}
$ cd build && ./parent
[Parent PID 1132] Started.
[Parent PID 1132] Enter output filename: result.txt
[Parent PID 1132] Enter numbers (space separated) or 'quit' to exit:
1.5 2.0 3.0
[Child PID 1133] Started. Reading from stdin, writing to 'result.txt'
[Child PID 1133] Processed line, found 3 numbers, sum = 6.5
10.1 20.2
[Child PID 1133] Processed line, found 2 numbers, sum = 30.3
quit
[Child PID 1133] Finished and exited.
[Parent PID 1132] Child (PID 1133) exited with status 0
\end{verbatim}

Строки с \texttt{PID 1132} и \texttt{PID 1133} подтверждают, что процессы выполняются конкурентно. Статус выхода 0 и корректные суммы (6.5 и 30.3) подтверждают правильность вычислений и отсутствие зомби-процессов.

\section{Выводы}
В ходе выполнения лабораторной работы была реализована многопроцессная программа на языке C++, демонстрирующая механизмы управления процессами и межпроцессного взаимодействия в POSIX-совместимой ОС. Родительский процесс создаёт дочерний через \texttt{fork()} и \texttt{execl()}, организует обмен данными через канал (\texttt{pipe}) с перенаправлением ввода посредством \texttt{dup2()}. Корректность работы подтверждена совпадением вычисленных сумм с ожидаемыми значениями и нормальным кодом завершения. Освоены практические навыки системного программирования и работы с POSIX API.

\end{document}